\documentclass[11pt]{article}

\usepackage{amsmath}
\usepackage{amssymb}
\usepackage{amsthm}
\usepackage{xurl}
\usepackage[colorlinks=true,linkcolor=blue,citecolor=blue,urlcolor=blue]{hyperref}
\setlength{\emergencystretch}{2em}

% ---------- theorem environments ----------
\theoremstyle{plain}
\newtheorem{theorem}{Theorem}[section]
\newtheorem{lemma}[theorem]{Lemma}
\newtheorem{corollary}[theorem]{Corollary}
\newtheorem{proposition}[theorem]{Proposition}
\newtheorem{claim}[theorem]{Claim}
\theoremstyle{definition}
\newtheorem{definition}[theorem]{Definition}
\newtheorem{conjecture}[theorem]{Conjecture}
\theoremstyle{remark}
\newtheorem{remark}[theorem]{Remark}

% ---------- operators ----------
\DeclareMathOperator{\supp}{supp}
\DeclareMathOperator{\ex}{ex}
\DeclareMathOperator{\mean}{mean}
\DeclareMathOperator{\tr}{tr}
\DeclareMathOperator{\Int}{int}

\newcommand{\PP}{\mathcal P}
\newcommand{\DD}{\Delta}
\newcommand{\EE}{\mathbb E}
\newcommand{\eps}{\varepsilon}
\newcommand{\rh}{\operatorname{rh}}
\newcommand{\sstar}{s^{\star}}
\newcommand{\Lb}{\mathrm{L}}
\newcommand{\Bb}{\mathrm{B}}
\newcommand{\Cc}{\mathrm{C}}
\newcommand{\Qq}{\mathrm{Q}}
\newcommand{\Ee}{\mathrm{E}}
\newcommand{\Ff}{\mathrm{F}}
\newcommand{\Lh}{\Lambda}
\newcommand{\MM}{\mathcal M}
\newcommand{\PROVED}{\textsc{proved}}
\newcommand{\CITED}{\textsc{cited}}
\newcommand{\NUMERICAL}{\textsc{numerical}}

\title{An explicit constant exceeding $(3-\sqrt5)/2$ for union-closed families:
a machine-checked certificate with independent replay}
\author{The union-closed certification project\thanks{Repository artifact:
\texttt{math/uc/} of the \texttt{jinleic-workspace} tree; every claim carries a
machine-checked artifact named in the text.}}
\date{\today}

\begin{document}

\maketitle

\begin{abstract}
Let
\[
\psi=\frac{3-\sqrt5}{2}
\quad\text{and}\quad
t_{\mathrm{cert}}
=\frac{955165028125263}{2500000000000000}
=0.3820660112501052.
\]
We establish the first explicit union-closed frequency constant strictly above
$\psi$ with a complete proof chain: the exact rational comparison
$t_{\mathrm{cert}}\ge\psi+10^{-4}$ and the required five-parameter functional
inequality are \PROVED{} by machine-checked artifacts, while the single
entropy-to-union-closed implication is \CITED{} from Cambie
\cite[Question~2 and \S4]{cambie}.  Thus every finite union-closed family
$\mathcal F\ne\{\emptyset\}$ has an element in at least
$t_{\mathrm{cert}}|\mathcal F|$ member sets, subject precisely to that cited
implication.  The new proof decomposes the OR-entropy kernel into one positive
square minus an explicit positive-semidefinite defect, reduces the global
minimum to two pair-orbits (five parameters), and certifies nonnegativity by an
Arb interval branch-and-bound.  A frozen collector independently replayed all
eight one-byte-per-node proof traces and printed
\texttt{COMPOSITE CERTIFICATE}.  As a second, independent contribution, we
prove Liu's previously only numerically verified positive-semidefiniteness
hypothesis behind the best published constant, $0.382709087918741$; his
separate nine-parameter global-optimization hypothesis remains open, so that
constant is still conditional.  We do not claim Frankl's constant $1/2$,
optimality of $t_{\mathrm{cert}}$, certification of the reported
$0.3823455\ldots$ or $0.3827090\ldots$ constants, or a reproof of the one cited
implication.
\end{abstract}

\section{Introduction}

\subsection{Frankl's conjecture and the pre-constant era}
A finite family of sets is \emph{union-closed} if $A\cup B$ belongs to the
family whenever $A$ and $B$ do.  The union-closed sets conjecture says that
every finite union-closed family $\mathcal F\ne\{\emptyset\}$ contains an
element in at least $|\mathcal F|/2$ member sets.  It is commonly attributed to
P\'eter Frankl in 1979, although the historical record describes it as a
folklore conjecture already known in the 1970s
\cite[\S1]{bruhn}.  We therefore use ``Frankl's conjecture'' as the conventional
name, not as a claim of sole priority.

Before the 2022 entropy breakthrough, the general quantitative bounds decayed
with the size of the family.  Knill \cite{knill} and W\'ojcik \cite{wojcik}
obtained frequency fractions of order $\Omega(1/\log|\mathcal F|)$; in
particular these were not positive constants independent of
$|\mathcal F|$.  The relevant early author is Emanuel \emph{Knill}.

\subsection{The entropy ladder: 1/100 and the golden-ratio barrier}
Gilmer \cite{gilmer} introduced the entropy method that produced the first
dimension-free constant, proving $1/100$.  For a regular Borel probability law
$\mu$ on $[0,1]$, its iid OR-entropy and one-coordinate entropy are
\[
\Qq(\mu)=\iint h(p+q-pq)\,d\mu(p)\,d\mu(q),
\qquad
\Lb(\mu)=\int h(p)\,d\mu(p),
\]
where $h(u)=-u\log_2u-(1-u)\log_2(1-u)$, continuously extended by
$h(0)=h(1)=0$.  The sharp iid inequality was then established independently
in near-simultaneous work of Alweiss--Huang--Sellke
\cite{ahs}, Chase--Lovett \cite{chaselovett}, Sawin \cite{sawin}, and Pebody
\cite{pebody}.  It yields
\[
\psi=\frac{3-\sqrt5}{2}\approx0.3819660112501051,
\]
the largest explicit constant above which no complete proof had previously
been available.  Chase--Lovett also showed that $\psi$ is sharp for their
approximate-union-closed formulation, and it is the sharp barrier for the iid
entropy strategy.

Sawin's revised paper \cite[\S2]{sawin} does more than prove the $\psi$ bound:
it sketches how a convex combination of the iid coupling and a dependent
max-entropy coupling gives some unspecified constant strictly larger than
$\psi$.  The sketch deliberately does not calculate an explicit improvement.
This is the conceptual second strategy used by Yu, Cambie, and the present
work.

\subsection{Claims beyond the golden-ratio barrier and their proof status}
The distinction among \PROVED{}, \CITED{}, and \NUMERICAL{} statements is
load-bearing here.

\paragraph{Yu.}
Yu \cite[\S4]{yu} claims the explicit value
$0.38234\ldots$ (arXiv:2212.00658v2).  The argument invokes
Krein--Milman/Carath\'eodory and then a concavity inequality to pass from
extreme points to a finite-support optimization.  That inference is false on
Yu's own extreme-point class.  The repository artifact
\texttt{uc/yu\_gap.py} gives a particularly simple structured counterexample:
with $Q_1=\delta_{(0.15,0.15)}$,
$Q_2=(1-\beta)\delta_{(0.10,0.10)}+\beta Q_{1,1}$ at mean
$0.382345533366702721$, and mixing weight $1/2$, it reports the claimed
concavity defect as $-2.476271\times10^{-2}$.  This displayed sign is a
\NUMERICAL{} diagnostic, not a proof step in our theorem; it refutes the
specific reduction with a margin far from roundoff.  Theorem~\ref{thm:B3}
below supplies a different support reduction with all hypotheses retained.

\paragraph{Cambie.}
Cambie \cite{cambie} formulates the two-strategy functional inequality as
Question~2 and reports
\[
c^\star\approx0.3823455333667027.
\]
Version~2 explicitly describes its finite-dimensional verification as follows:
\emph{``In contrast to the detailed work in e.g.\ [1] for the constant
$(3-\sqrt5)/2$, this is done slightly less rigorous.''}
In other words, that optimization was carried out slightly less rigorously:
the paper itself notes local-minimum and finite-precision issues.  We do not
use that numerical verification.  We use only the separate implication in
Cambie's \S4 that converts a valid answer to Question~2 into the
union-closed frequency bound; that one imported step is always labeled
\CITED{} below.

\paragraph{Liu.}
Liu \cite{liu} develops conditionally iid couplings.  The explicit constant in
Theorem~13,
\[
0.382709087918741,
\]
is conditional on two hypotheses stated there: the positive-semidefiniteness
hypothesis of \S V-A and the global-minimizer-structure hypothesis of
\S V-B.  Both were reported as numerically verified rather than proved.  As a
second, independent contribution, Theorem~\ref{thm:liu-psd} proves the
\S V-A hypothesis.  The nine-parameter \S V-B hypothesis remains open, so
Liu's decimal is still conditional and is not a theorem claimed here.
Related entropy developments include Wakhare \cite{wakhare} and
Lu--Raz \cite{luraz}.

\paragraph{Quantitative ladder.}
\begin{itemize}
\item \PROVED{} (size-dependent):
$\Omega(1/\log|\mathcal F|)$, by Knill \cite{knill} and W\'ojcik
\cite{wojcik}.
\item \PROVED{}:
$1/100$, by Gilmer \cite{gilmer}.
\item \PROVED{}:
$\psi=(3-\sqrt5)/2$, by Alweiss--Huang--Sellke \cite{ahs},
Chase--Lovett \cite{chaselovett}, Sawin \cite{sawin}, and Pebody
\cite{pebody}.
\item Claimed with a finite-dimensional proof gap:
$c^\star\approx0.3823455333667027$, by Yu \cite{yu} and Cambie
\cite{cambie}.
\item Conditional because the nine-parameter Hypothesis~2 remains open
(Hypothesis~1 is proved in this work):
$0.382709087918741$, in Liu's Theorem~13 \cite{liu}.
\item \PROVED{} functional inequality plus one \CITED{} implication:
$t_{\mathrm{cert}}=0.3820660112501052$, in this work and Cambie's
\S4 \cite{cambie}.
\end{itemize}
The last item is the exact rational claim
$t_{\mathrm{cert}}\ge\psi+10^{-4}$, never an equality between the rational
target and the irrational number $\psi+10^{-4}$.

\subsection{The result proved here}
For a compact metric space $X$, let $\MM(X)$ be the real vector space of finite
signed regular Borel measures, $\MM_+(X)$ its positive cone, and
\[
\PP(X)=\{\mu\in\MM_+(X):\mu(X)=1\}.
\]
All measures and couplings in the theorem statements are regular Borel
measures.  For $\mu\in\PP([0,1])$ put
$\mean(\mu)=\int_0^1p\,d\mu(p)$ and
$\Bb(\mu)=1-\mean(\mu)$.  A self-coupling of $\mu$ means an element of
$\PP([0,1]^2)$ whose two marginals both equal $\mu$; the set of all such
couplings is denoted $\Pi(\mu,\mu)$.  Weak-$*$ topology is used only where it
is named explicitly.

Set
\[
\begin{aligned}
\alpha_0&=\frac{356069}{10000000}=0.0356069,\\
t_{\mathrm{cert}}
&=\frac{955165028125263}{2500000000000000}\\
&=0.3820660112501052.
\end{aligned}
\]
The rational inequality
$t_{\mathrm{cert}}\ge\psi+10^{-4}$ is checked exactly by the campaign
preflight.

\begin{theorem}[Machine-checked functional inequality]
\label{thm:main}
Let $\alpha_0=0.0356069$ and let
$t_{\mathrm{cert}}=0.3820660112501052$ denote the exact rationals displayed
above.  For every $\mu\in\PP([0,1])$ satisfying
$\mean(\mu)\le t_{\mathrm{cert}}$, define
\begin{align*}
\Lb(\mu)&=\int h(p)\,d\mu(p),\\
\Qq(\mu)&=\iint h(p+q-pq)\,d\mu(p)\,d\mu(q),\\
\Cc(\mu)&=\min_{\pi\in\Pi(\mu,\mu)}
  \iint h\bigl(\sstar(p,r)\bigr)\,d\pi(p,r),\\
\sstar(p,r)&=\max\bigl(p,r,\min(p+r,\tfrac12)\bigr).
\end{align*}
Then
\begin{equation}
\Ff(\mu):=(1-\alpha_0)\Qq(\mu)+\alpha_0\Cc(\mu)-\Lb(\mu)\ge0.
\label{eq:Fdef}
\end{equation}
\end{theorem}

\begin{corollary}[Union-closed consequence; one cited implication]
\label{cor:uc}
Assume the entropy-to-union-closed implication in Cambie's Question~2 and
\S4 \cite{cambie}: for $\alpha\in[0,1]$ and $t\in[0,1]$, validity of the
corresponding functional inequality for every
$\mu\in\PP([0,1])$ with $\mean(\mu)\le t$ implies the frequency bound $t$
for every finite union-closed family other than $\{\emptyset\}$.  Then every
finite union-closed family $\mathcal F\ne\{\emptyset\}$ has an element
belonging to at least
\[
t_{\mathrm{cert}}|\mathcal F|
=0.3820660112501052\,|\mathcal F|
\]
member sets.
\end{corollary}

Theorem~\ref{thm:main} is \PROVED{} here.  The implication explicitly assumed
in Corollary~\ref{cor:uc} is \CITED{} and is not re-proved.  Together they give
the first explicit constant strictly above $\psi$ with a complete proof chain.
The scope is exact: we do not prove Frankl's $1/2$, certify $c^\star$, certify
Liu's conditional decimal, prove optimality, or claim a result at any larger
target.

The proof has four structural ingredients:
\begin{enumerate}
\item[(A)] an exact OR-entropy-kernel decomposition; see
\S\ref{sec:scaffold} and \path{decomposition.py};
\item[(A${}'$)] the identity $\Qq=2\Bb\Lb-\Ee$ with an explicit
positive-semidefinite defect; see \S\ref{sec:scaffold} and
\path{reduction.py};
\item[(B${}'''$)] a support reduction to at most two pair-orbits, hence five
parameters; see \S\ref{sec:B3} and \path{thmB3_proof.py};
\item a scale-free margin lemma and the $\rh$ bound controlling the entropy
sinks; see \S\ref{sec:margin}, \path{margin_lemma.py}, and
\path{lemma_rh_proof.py}.
\end{enumerate}
An Arb interval branch-and-bound then certifies the reduced functional
(\S\ref{sec:certificate}).  Campaign I committed all eight slices, totaling
$488{,}465{,}854$ boxes with zero residual and empty final stacks.  Its frozen
collector replayed every trace and printed \texttt{COMPOSITE CERTIFICATE} on
2026-08-21.  Appendix~\ref{sec:verify} gives the exact replay command and
trust base.  Every search heuristic or sampled diagnostic is labeled
\NUMERICAL{} and is not a proof step.

\section{The exact scaffold: Theorems A and A-prime}
\label{sec:scaffold}

Throughout, logarithms are natural unless the subscript $2$ shows an explicit base.
We use the natural-log kernels
\begin{equation}
H(u)=-u\ln u-(1-u)\ln(1-u),\qquad
T(u)=u+(1-u)\ln(1-u),
\label{eq:Hdef}
\end{equation}
extended continuously by $H(0)=0$, $T(0)=0$, $T(1)=1$; recall that
$h(u)=H(u)/\ln2$.

\begin{lemma}[T-series]
\label{lem:Tseries}
For $0\le u\le1$,
\begin{equation}
T(u)=\sum_{n=2}^{\infty}\frac{u^n}{n(n-1)},\qquad
T(1)=\sum_{n=2}^{\infty}\frac{1}{n(n-1)}=1.
\label{eq:Tseries}
\end{equation}
\end{lemma}

\begin{proof}
$T'(u)=-\ln(1-u)=\sum_{m\ge1}u^{m}/m$ on $[0,1)$, so integrating from $0$
gives $T(u)=\sum_{m\ge1}u^{m+1}/(m(m+1))=\sum_{n\ge2}u^n/(n(n-1))$; the
endpoint follows by Abel's theorem or by the telescoping $1/(n(n-1))=1/(n-1)-
1/n$. The machine artifact \texttt{decomposition.py} checks the closed
form against the series at five points to $10^{-30}$ and the telescoping sum
exactly. \qedhere
\end{proof}

\begin{theorem}[Theorem A; \texttt{decomposition.py}]
\label{thm:A}
Let $H,T$ be the continuous kernels in \eqref{eq:Hdef}, and set
$g(x)=x(1-\ln x)$ and $a(x)=-x\ln x$, with $g(0)=a(0)=0$.
For every $x,y\in[0,1]$,
\begin{equation}
H(xy)=g(x)\,g(y)-a(x)\,a(y)-T(xy).
\label{eq:kernel}
\end{equation}
Moreover $T(xy)=\sum_{n\ge2}(xy)^n/(n(n-1))$ is positive semidefinite in
the following explicit measure-cone sense:
\[
\iint T(xy)\,d\eta(x)\,d\eta(y)\ge0
\qquad\text{for every }\eta\in\MM([0,1]).
\]
Thus \eqref{eq:kernel} writes the OR-entropy kernel as the single positive
square $g\otimes g$ minus the positive-semidefinite remainder
$a\otimes a+T$.
\end{theorem}

\begin{proof}
Since $g=a+\mathrm{id}$,
\[
g(x)g(y)-a(x)a(y)=x\,a(y)+y\,a(x)+xy=-xy\ln(xy)+xy
=H(xy)+T(xy),
\]
the last step by $H(u)+T(u)=-u\ln u+u$ from \eqref{eq:Hdef}. Rearranging
gives \eqref{eq:kernel}.  The \PROVED{} gate in
\texttt{decomposition.py} asserts the identity as an exact SymPy tautology on
three symbolic atoms.  Its checks on 200 random signed measures and on a
$200\times200$ grid are additional \NUMERICAL{} controls, not proof steps.
\end{proof}

\begin{remark}
The ``at most one positive square'' statement is the algebraic heart of the method:
it says the two-strategy entropy $\Qq$ differs from a square of a linear
functional by a positive-semidefinite remainder, so that a variance
(``defect'') term appears with a sign we can control. This is the property that
makes the later support reduction possible; see \S3.
\end{remark}

We now pass to the probability-theoretic formulation.  For
$\mu\in\PP([0,1])$ write, with the substitution $x=1-p$ in the $W$-kernel,
\begin{equation}
\begin{split}
\Ee(\mu)&=\frac1{\ln2}\iint W(1-p,\,1-q)\,d\mu(p)\,d\mu(q),\\
W(x,y)&=xy-x\,T(y)-y\,T(x)+T(xy).
\end{split}
\label{eq:Wdef}
\end{equation}

\begin{theorem}[Theorem A${}'$; \texttt{reduction.py}]
\label{thm:Ap}
Let $W$ and $\Ee$ be as in \eqref{eq:Wdef}.  For every $x,y\in[0,1]$,
\begin{equation}
W(x,y)=\sum_{n=2}^{\infty}\frac{(x-x^n)(y-y^n)}{n(n-1)}.
\label{eq:Wseries}
\end{equation}
Consequently $W$ is positive semidefinite on the full signed-measure space:
\[
\iint W(x,y)\,d\eta(x)\,d\eta(y)\ge0
\qquad\text{for every }\eta\in\MM([0,1]).
\]
For every $\mu\in\PP([0,1])$, with
$\widetilde\mu=(p\mapsto1-p)_\#\mu$ and
$M_n=\int x^n\,d\widetilde\mu(x)$, one has
\begin{align}
\Qq(\mu) &= 2\,\Bb(\mu)\,\Lb(\mu)-\Ee(\mu), \label{eq:Q2BL}\\
\ln2\cdot\Ee(\mu) &= \sum_{n=2}^{\infty}\frac{(M_1-M_n)^2}{n(n-1)},
  \label{eq:Emoment}\\
\frac{W(x,x)}{\ln2} &= 2x\,h(x)-h(x^2)\qquad(0\le x\le1).
\label{eq:Wdiag}
\end{align}
In particular $\Ee(\mu)\ge0$ for every $\mu\in\PP([0,1])$.
\end{theorem}

\begin{proof}
For \eqref{eq:Wseries}, expand termwise over the four convergent series of
\eqref{eq:Wdef} using \eqref{eq:Tseries}: the telescoping identity
$\sum_{n\ge2}1/(n(n-1))=1$ and the two evaluations
$\sum_{n\ge2}x y^n/(n(n-1))=x\,T(y)$ and its symmetric partner reproduce
$W$. Each summand $(x-x^n)(y-y^n)$ is a rank-one kernel, so $W$ is a sum
of nonnegative squares; its quadratic form against a finite signed measure is
therefore a sum of squares, giving \eqref{eq:Emoment} with
$\int(x-x^n)\,d\widetilde\mu=M_1-M_n$.

For \eqref{eq:Q2BL}, note that $p+q-pq=1-xy$ with $x=1-p$, $y=1-q$,
so $h(p+q-pq)=H(xy)/\ln2$, and by Theorem~\ref{thm:A}
\[
\ln2\cdot\Qq = \Bigl(\int g\,d\widetilde\mu\Bigr)^2
-\Bigl(\int a\,d\widetilde\mu\Bigr)^2-\iint T(xy)\,d\widetilde\mu\,
d\widetilde\mu .
\]
Let $\overline T=\int T\,d\widetilde\mu$ and $B=\int x\,d\widetilde\mu$;
then $\int g\,d\widetilde\mu=\int a\,d\widetilde\mu+B$, and from
$H(x)=a(x)+x-T(x)$,
\[
\ln2\cdot\Lb=\int a\,d\widetilde\mu+B-\overline T .
\]
A direct expansion gives
\[
\ln2\,(2B\,\Lb-\Qq)
=B^2-2B\,\overline T+\iint T(xy)\,d\widetilde\mu\,d\widetilde\mu
=\iint W\,d\widetilde\mu\,d\widetilde\mu=\ln2\cdot\Ee,
\]
which is \eqref{eq:Q2BL}. Finally
\eqref{eq:Wdiag} follows from the identity
$T(u)=-u\ln u+u-H(u)$:
$W(x,x)=x^2-2x\,T(x)+T(x^2)=2x\,H(x)-H(x^2)$.
The \PROVED{} gate in \texttt{reduction.py} verifies
\eqref{eq:Wseries} against the closed form within an analytic tail bracket
(the remainder is at most $1/N$ after $N$ terms) and asserts the symbolic
identity chain.  Its random-law evaluations and four-point evaluations,
including the golden point, are \NUMERICAL{} controls only.
\end{proof}

Substituting \eqref{eq:Q2BL} into \eqref{eq:Fdef} yields the form used
throughout:

\begin{corollary}[Defect form of the functional]
\label{cor:Fid}
For every $\alpha\in[0,1]$ and every $\mu\in\PP([0,1])$, with
$\Ff(\mu)=(1-\alpha)\Qq(\mu)+\alpha\Cc(\mu)-\Lb(\mu)$,
\begin{equation}
\Ff(\mu)=\bigl[2(1-\alpha)\,\Bb(\mu)-1\bigr]\,\Lb(\mu)
+\alpha\,\Cc(\mu)-(1-\alpha)\,\Ee(\mu).
\label{eq:Fid}
\end{equation}
\end{corollary}

\subsection{The iid golden point and the golden-ratio barrier}
For $\alpha\in[0,1]$ and a point mass $\mu=\delta_p$ with $p\in[0,1]$,
the only self-coupling is $\delta_{(p,p)}$, and
\begin{equation}
\Ff(\delta_p)=\begin{cases}
h\bigl((1-p)^2\bigr)-h(p),&\alpha=0,\\
(1-\alpha)\,h\bigl((1-p)^2\bigr)+\alpha\,h\bigl(\sstar(p,p)\bigr)-h(p),
&\alpha\ge0,
\end{cases}
\label{eq:Fpoint}
\end{equation}
since $\Qq(\delta_p)=h(2p-p^2)=h(1-(1-p)^2)$ and $\Cc(\delta_p)=
h(\sstar(p,p))$.

\begin{corollary}[The iid point-mass barrier; \texttt{reduction.py}]
\label{cor:psi}
For every $p\in[0,1]$, at $\alpha=0$,
\[
\Ff(\delta_p)\ge0
\iff (1-p)^2\ge p
\iff p\le\psi=\frac{3-\sqrt5}{2}.
\]
\end{corollary}

For $p\in[0,\tfrac12]$, both $(1-p)^2$ and $p$ lie in
$[0,\tfrac12]$ whenever the inequality can change sign, and $h$ is increasing
there; for $p>\tfrac12$, one has
$h((1-p)^2)<h(1-p)=h(p)$.  The quadratic inequality has smaller root
$\psi$.  At the golden point $x=1/\varphi$, $x^2=1-x$ exactly, so
$\Ff$ vanishes: $\psi$ is precisely where the PSD defect $\Ee$ first
overtakes $(2\Bb-1)\Lb$.

For later comparisons, define the two-orbit obstruction constants without a
decimal approximation.  Let $b^\star$ be the larger root in $(0,\tfrac12)$ of
\[
h(b)\,[2-h(b)]=h\bigl((1-b)^2\bigr),
\]
put $a^\star=[1-h(b^\star)]/[2-h(b^\star)]$, and set
$c^\star=a^\star+(1-a^\star)b^\star$.

The endpoint value
$\Ff(\delta_\psi)=\alpha[1-h(\psi)]\approx0.0405813\,\alpha>0$
for $\alpha>0$, so point masses do not cap the constant once the second
strategy is active.  The crossover value is
\begin{equation}
\alpha_{\mathrm{crit}}=
\frac{h(c^{\star})-h\bigl((1-c^{\star})^2\bigr)}
{1-h\bigl((1-c^{\star})^2\bigr)}
\approx0.0144054585140081.
\label{eq:acrit}
\end{equation}

\begin{lemma}[Point-mass threshold; \texttt{reduction.py} \S6]
\label{lem:acrit}
Let $\alpha\in[0,1)$, let
$\Ff(\delta_p)=(1-\alpha)h(2p-p^2)
+\alpha h(\sstar(p,p))-h(p)$ for $p\in[\tfrac14,\tfrac12]$, and let
$c^\star,\alpha_{\mathrm{crit}}$ be defined above.  The equation
$\Ff(\delta_p)=0$ has a unique root in $(\tfrac14,\tfrac12)$, and
\[
\text{that root exceeds }c^{\star}
\iff \alpha>\alpha_{\mathrm{crit}} .
\]
For $\alpha=\alpha_0=0.0356069>\alpha_{\mathrm{crit}}$, point masses are
therefore not binding.
\end{lemma}

\begin{proof}[Proof sketch; complete in \texttt{reduction.py}]
On $(1/4,1/2)$, $\sstar(p,p)=1/2$ is \emph{constant}
(since $p\le1/2\le p+p$), so the $\alpha$-term drops out of
$\partial_p\Ff(\delta_p)$, giving
\[
\frac{d}{dp}\,\Ff(\delta_p)=(1-\alpha)(2-2p)\,h'\bigl(2p-p^2\bigr)-h'(p).
\]
On $[p_0,\frac12)$ with $p_0=1-1/\sqrt2$ where $2p-p^2\ge\frac12$, each
factor has a proved sign: $h'(2p-p^2)\le0$, $2-2p>0$, $h'(p)>0$,
so the derivative is strictly negative. On $[0,p_0]$, both $h(2p-p^2)\ge
h(p)$ and $h(\sstar(p,p))\ge h(p)$ hold uniformly, so $\Ff\ge0$ for
\emph{every} $\alpha\in[0,1]$; and
$\Ff(\delta_{1/2})<0$ for $\alpha<1$.  Hence there is a unique root.  Its
position relative to $c^{\star}$ is decided by the sign at $c^{\star}$,
which is affine in $\alpha$ and vanishes exactly at \eqref{eq:acrit}.  The
artifact asserts this identity chain.  Its bisection over a grid of 401 points
per $\alpha$ is an additional \NUMERICAL{} control, not the assertion.
\end{proof}

\subsection{Why the variance term is load-bearing}
Dropping the variance in the crudest way --- replacing $\Ee$ by a pointwise
upper bound derived from \eqref{eq:Wdiag} --- gives a bound whose ratio form is
governed by $\kappa^{\star}=\max_{x\in(0,1)}[2x-h(x^2)/h(x)]$, i.e.\ an
explicit constant $(1-\kappa^{\star})/2\approx0.34753<\psi$ (asserted in
\texttt{reduction.py} \S8). The maximizer sits at $p^{\star}\approx0.117$,
far below the mean bound, so the pointwise route \emph{wastes} the constraint
$\EE p\le t$: the variance term is load-bearing, and one must either retain it
(as we do, via the support reduction and the certificate) or control it sharply
enough to win it back.

\section{Support reduction: Theorem B-triple-prime}
\label{sec:B3}

The defect term $\Ee$ is convex in $\mu$ (a positive semidefinite quadratic
form), so the certificate cannot simply minimize a concave relaxation; instead we prove
that the global minimum of the exact functional $\Ff$ is attained on a tiny
parameter family. The functional must first be recast on \emph{pair-orbit
measures}.

\begin{definition}[Pair-orbit measures]
\label{def:pair}
Let
\[
\DD=\{(p,q)\in[0,1]^2:p\le q\},
\]
and let $\PP(\DD)$ denote the regular Borel probability measures on this
compact triangle.  A point $(p,q)\in\DD$ is a \emph{pair orbit}.  Every
$\nu\in\PP(\DD)$ induces the marginal law
\begin{align*}
\mu_\nu=M(\nu)
&=\frac{(\pi_1)_\#\nu+(\pi_2)_\#\nu}{2}\in\PP([0,1]),\\
\int f\,d\mu_\nu
&=\int_{\DD}\frac{f(p)+f(q)}2\,d\nu(p,q)
\qquad(f\in C[0,1]).
\end{align*}
A \emph{two-orbit law} is
\[
\nu=w\delta_{(p_1,q_1)}+(1-w)\delta_{(p_2,q_2)},
\quad
0\le p_i\le q_i\le1,\quad 0\le w\le1.
\]
It is parametrized by $(p_1,q_1,p_2,q_2,w)$.
\end{definition}

The point of the folding is exactness of the coupling cost. If $\pi$ is a
symmetric coupling of $\mu$ with itself (which suffices: replacing $\pi$ by
$(\pi+\pi^{\top})/2$ preserves both marginals and every symmetric cost), folding
$(p,r)\mapsto(\min(p,r),\max(p,r))$ yields a law $\nu$ on $\DD$ whose
induced marginal is $\mu$ and whose $\sstar$-cost is unchanged. Hence
\[
\Cc(\mu)=\min_{\substack{\nu\in\PP(\DD)\\\mu_\nu=\mu}}\int c\,d\nu,\qquad
c(p,q)=h\bigl(\sstar(p,q)\bigr),
\]
and the joint optimization over laws and couplings collapses to
\[
\inf_{\substack{\mu\in\PP([0,1])\\\mean(\mu)\le t}}\Ff(\mu)
=
\inf_{\substack{\nu\in\PP(\DD)\\\mean(\mu_\nu)\le t}}\Phi_\alpha(\nu),
\]
where $\Phi_\alpha$ is the functional stated explicitly in
Theorem~\ref{thm:B3}.

\begin{theorem}[Theorem B${}'''$; \texttt{thmB3\_proof.py}]
\label{thm:B3}
Let $\DD$ and the map $\nu\mapsto\mu_\nu$ be as in
Definition~\ref{def:pair}.  For every $\alpha\in[0,1]$ and $t\in[0,1]$,
define on the weak-$*$ compact convex set $\PP(\DD)$
\[
\Phi_\alpha(\nu)
=(1-\alpha)\Qq(\mu_\nu)
+\alpha\int_{\DD}h(\sstar(p,q))\,d\nu(p,q)-\Lb(\mu_\nu).
\]
The minimum of $\Phi_\alpha$ over
$\{\nu\in\PP(\DD):\mean(\mu_\nu)\le t\}$ is attained by a measure supported
on at most two points of $\DD$.  Its induced marginal is therefore supported
on at most four points of $[0,1]$.
\end{theorem}

The proof uses the \CITED{} Riesz, Banach--Alaoglu, Bauer, and moment-set
theorems identified below.  The complete argument carries their hypotheses
explicitly.  The executable \texttt{thmB3\_proof.py} also runs sampled
\NUMERICAL{} controls, none of which is a proof step.

\begin{proof}
\emph{Step 1: compactness.} $\DD$ is a closed subset of $[0,1]^2$, hence a
compact metric space. By the Riesz representation theorem (van~Neerven
\cite{vanneerven}, Thm.~4.2), regular Borel measures on $\DD$ are identified
with $C(\DD)^*$, and by Banach--Alaoglu (van~Neerven \cite{vanneerven},
Thm.~4.50) the closed unit ball of $C(\DD)^*$ is weak-$*$ compact. The set
$\PP(\DD)$ is weak-$*$ closed there (it is the intersection of closed conditions
$\langle 1,\nu\rangle=1$ and $\langle f,\nu\rangle\ge0$ for all $f\ge0$),
so $\PP(\DD)$ is weak-$*$ compact. (Equivalently, Prokhorov: every family of
laws on compact $\DD$ is tight.)

\emph{Step 2: continuity.} The functions
\begin{align*}
g(p,q)&=\frac{2-p-q}{2},\\
\ell(p,q)&=\frac{h(p)+h(q)}{2},\\
\sstar(p,q)&=\min\bigl(\max(\tfrac12,p,q),\min(p+q,1)\bigr).
\end{align*}
and $c=h\circ\sstar$
are continuous on $\DD$ ($r\ln r\to0$ gives continuity of $h$ at the
endpoints; min and max preserve continuity). The kernel $W$ of \eqref{eq:Wdef}
is continuous on $[0,1]^2$ and $W(0,\cdot)=W(1,\cdot)=0$ by direct
substitution, so $T$ and $W$ are likewise continuous. The marginal map $M$ is
continuous and linear. If $\mu_i\to\mu$ weak-$*$, then $\mu_i\otimes\mu_i\to
\mu\otimes\mu$ weak-$*$: product functions $a(x)b(y)$ separate points and form a
unital subalgebra, hence are uniformly dense in $C([0,1]^2)$
(Stone--Weierstrass; van~Neerven \cite{vanneerven}, Thm.~2.5), and
uniform approximation extends convergence to the continuous kernel $W(1-x,1-y)$.
Thus $\nu\mapsto\Ee(\mu_\nu)$ is weak-$*$ continuous, and so is
$\Phi_\alpha$ by \eqref{eq:Fid}.

\emph{Step 3: global attainment.}
The feasible set
$H_t=\{\nu\in\PP(\DD):\mean(\mu_\nu)\le t\}
=\{\nu\in\PP(\DD):\Bb(\mu_\nu)\ge1-t\}$ is the preimage of a closed ray
under a continuous map, hence weak-$*$ compact, and contains
$\delta_{(0,0)}$.  A global minimizer $\nu^{\star}$ exists; put
$\beta^{\star}=\Bb(\mu_{\nu^{\star}})$.

\emph{Step 4: the slice.}
The slice
$K_\beta=\{\nu\in\PP(\DD):\Bb(\mu_\nu)=\beta^{\star}\}$
is nonempty, convex, weak-$*$ compact, and contained in $H_t$, so
\[
\min_{K_\beta}\Phi_\alpha
=\Phi_\alpha(\nu^{\star})
=\min_{H_t}\Phi_\alpha.
\]
On $K_\beta$ the marginal mean, hence $\Bb$, is fixed.  By
\eqref{eq:Q2BL},
$\Qq(\mu_\nu)=2\beta^\star\Lb(\mu_\nu)-\Ee(\mu_\nu)$, and
\[
\Phi_\alpha(\nu)
=\bigl[2(1-\alpha)\beta^\star-1\bigr]\Lb(\mu_\nu)
+\alpha\int_{\DD}c\,d\nu-(1-\alpha)\Ee(\mu_\nu).
\]
The first two terms are weak-$*$ continuous linear functionals. For
$0\le\theta\le1$ and $\nu_0,\nu_1\in K_\beta$, positive semidefiniteness of
$W$ gives
\[
\Ee\bigl(\theta\mu_{\nu_0}+(1-\theta)\mu_{\nu_1}\bigr)
\le\theta\,\Ee(\mu_{\nu_0})+(1-\theta)\,\Ee(\mu_{\nu_1}),
\]
so $-(1-\alpha)\Ee$ is concave (here the restriction $\alpha\le1$ is used
once). Hence $\Phi_\alpha$ is continuous and concave on the compact convex
$K_\beta$.

\emph{Step 5: Bauer's minimum principle.}
By the Bauer maximum principle (Bauer \cite{bauer}; the formulation used here is
Bru--de~Siqueira~Pedra \cite{brupedra}, Lemma~3.3, with the reference to
Bauer's original paper checked against the publisher record), a continuous convex
function on a compact convex set attains its maximum at an extreme point; applied
to $-\Phi_\alpha$ on $K_\beta$, there is an extreme point $\nu^{**}\in K_\beta$
with $\Phi_\alpha(\nu^{**})=\min_{K_\beta}\Phi_\alpha$.

\emph{Step 6: extreme points of a one-moment slice.}
$K_\beta=\{\nu:\int g\,d\nu=\beta^{\star}\}$ is a moment set cut from
$\PP(\DD)$ by \emph{one} non-constant moment condition (plus the mass
normalization, already part of $\PP(\DD)$). By Winkler's theorem on extreme
points of moment sets \cite{winkler}, as reproduced in Pinelis'
Theorem~12 \cite{pinelis}, every extreme point of such a slice is supported on at
most $1+1=2$ points. For completeness we give the short direct proof in this
case. Suppose an extreme $\nu\in K_\beta$ has three distinct support points, and
choose pairwise disjoint Borel neighborhoods $A_1,A_2,A_3$ of them with
positive $\nu$-mass. The three vectors
\[
v_i=\Bigl(\nu(A_i),\;\int_{A_i}g\,d\nu\Bigr)\in\mathbb R^2
\]
are linearly dependent: let $(a_1,a_2,a_3)\ne0$ with $\sum_i a_i v_i=0$,
and put $\eta=\sum_i a_i\,\nu|_{A_i}$. Then $\eta$ is a nonzero signed
measure with total mass $0$ and $g$-moment $0$, so for small $\eps>0$, the
two measures $\nu+\eps\eta$ and $\nu-\eps\eta$ (with densities $1\pm\eps a_i$
on $A_i$) are distinct nonnegative members of $K_\beta$ whose midpoint is $\nu$,
contradicting extremality. Hence $|\supp(\nu^{**})|\le2$.

\emph{Step 7.}
$\nu^{**}$ has $\Bb(\mu_{\nu^{**}})=\beta^{\star}$, so
$\mean(\mu_{\nu^{**}})=1-\beta^{\star}\le t$, and by Step~4 it is a global
minimizer.  Its marginal is supported on at most four points, and
$\nu^{**}$ is a two-orbit (or one-orbit) law in the five-parameter family of
Definition~\ref{def:pair}. \qedhere
\end{proof}

\begin{remark}[Why two orbits, and the controls in \texttt{thmB3\_proof.py}]
The decisive move is freezing $\Bb$ (equivalently the mean) up front: by
Theorem~A${}'$ that is the \emph{only} place the mean enters $\Qq$, so the
bilinear term becomes linear and the entire functional concave, leaving two moment
conditions (mass and $\Bb$) and hence two orbits.  The executable additionally
runs three explicitly labeled \NUMERICAL{} controls: fixed-$\Bb$ chord tests
(worst observed
$\Phi(\text{mid})-\text{chord}=+2.9\times10^{-9}$ over 3000 random
segments), a merged-atom continuity check of $M$ and $\Ee$, and a control that
\emph{removes} the fixed-$\Bb$ premise and detects non-concavity
($-0.124$ observed).
The binding obstruction (Section~\ref{sec:margin}) uses exactly two orbits, so
the theorem is tight in orbit count.
\end{remark}

\begin{corollary}[Five-parameter certification target]
\label{cor:5param}
Let $\alpha_0=0.0356069$ and
$t_{\mathrm{cert}}=0.3820660112501052$ be the exact rationals in
Theorem~\ref{thm:main}.  Then \eqref{eq:Fdef} holds for every
$\mu\in\PP([0,1])$ with $\mean(\mu)\le t_{\mathrm{cert}}$ if and only if
\[
\Phi_{\alpha_0}(\nu)\ge0
\]
for every
\[
\nu=w\delta_{(p_1,q_1)}+(1-w)\delta_{(p_2,q_2)},\qquad
0\le p_i\le q_i\le1,\quad 0\le w\le1,
\]
satisfying $\mean(\mu_\nu)\le t_{\mathrm{cert}}$.
\end{corollary}

\section{The margin lemma and the diagonal inequality}
\label{sec:margin}

The functional $\Phi_\alpha$ vanishes identically on the ``sink'' laws
$\{(1-u)\delta_0+u\delta_1\}$; on the way there both $\Lb$ and $\Ff$
collapse, so a raw sign test cannot terminate. The margin lemma removes this
degeneracy by dividing by $\Lb$.

For $x\in[0,1]$ define the diagonal density
\begin{align}
\rh(x)
&=\frac{W(1-x,1-x)}{\ln2}
=2(1-x)h(x)-h\bigl((1-x)^2\bigr),
  \label{eq:rhdiag}\\
\rho(x)&=\frac{\rh(x)}{h(x)}\qquad(0<x<1).
\end{align}
The positive-semidefinite series in Theorem~\ref{thm:Ap} gives a Hilbert-space
feature map $\mathbf v:[0,1]\to\ell^2$ satisfying
\begin{align*}
\langle\mathbf v(p),\mathbf v(q)\rangle
&=\frac{W(1-p,1-q)}{\ln2},\\
\|\mathbf v(p)\|^2&=\rh(p),\\
\Ee(\mu)&=\left\|\int\mathbf v\,d\mu\right\|^2.
\end{align*}

\begin{lemma}[Margin lemma; \texttt{margin\_lemma.py}]
\label{lem:margin}
Let $\alpha\in[0,1]$ and let $\mu\in\PP([0,1])$ satisfy
$\Lb(\mu)=\int h\,d\mu>0$.  With $\Cc(\mu)$ defined by the self-coupling
minimum in Theorem~\ref{thm:main}, put
\[
\sigma(\mu)=\int\sqrt{\rh(p)}\,d\mu(p),\qquad
\Lh(\mu)=2(1-\alpha)\Bb(\mu)-1
+\alpha\frac{\Cc(\mu)}{\Lb(\mu)}
-(1-\alpha)\frac{\sigma(\mu)^2}{\Lb(\mu)}.
\]
Then, for
$\Ff(\mu)=(1-\alpha)\Qq(\mu)+\alpha\Cc(\mu)-\Lb(\mu)$,
\begin{equation}
\Ff(\mu)\ge\Lb(\mu)\Lh(\mu).
\label{eq:margineq}
\end{equation}
Equality holds whenever
$|\supp(\mu)\setminus\{0,1\}|\le1$.
\end{lemma}

\begin{proof}
By the triangle inequality for Bochner integrals,
\[
\Ee(\mu)=\left\|\int\mathbf v\,d\mu\right\|^2
\le\left(\int\|\mathbf v\|\,d\mu\right)^2=\sigma(\mu)^2.
\]
Substitution in \eqref{eq:Fid} gives \eqref{eq:margineq}.  If the support has
at most one non-sink point, all nonzero feature vectors are collinear, so
equality holds. \qedhere
\end{proof}

Two properties make \eqref{eq:margineq} usable.
\begin{enumerate}
\item \emph{Bounded, never singular.}
The global Arb cover in \texttt{cert2.\_rho\_global} \PROVED{} the cap used
by the certificate.  Its approximate maximizer
$p^\star\approx0.1165$ and value
$\rho_{\max}\approx0.3049467$ are \NUMERICAL{} summaries.  Since
$\rh(x)\le\rho_{\max}h(x)$,
\[
\sigma(\mu)^2\le \rho_{\max}\Lb(\mu),
\]
so $\Lh$ is bounded and $\sigma^2/\Lb$ cannot diverge.  Along a sequence of
fixed mean $t$ collapsing to a sink, $\sigma^2$ is quadratic in the non-sink
mass while $\Lb$ is linear, and hence
\[
\Lh\longrightarrow 2(1-\alpha)(1-t)-1=\kappa_t
\approx0.19186255>0
\]
at the certified parameters.
\item \emph{Tight at the obstruction.}
With the exact $a^\star,b^\star$ defined above, the law
$\mu^\star=a^\star\delta_1+(1-a^\star)\delta_{b^\star}$ has exactly one
non-sink support point, so
\[
\Lh(\mu^\star)=\frac{\Ff(\mu^\star)}{\Lb(\mu^\star)}
\]
exactly.  The decimal summaries
$a^\star\approx0.0788772927059232$ and
$b^\star\approx0.3294547385030370$, the $10^{-12}$ evaluation, and the
Jensen-variant value $-1.94\times10^{-2}$ are \NUMERICAL{} diagnostics
recorded in \texttt{margin\_lemma.py}; none is a proof premise.
\end{enumerate}

The margin lemma is the reason the certificate can cover the near-total-sink
slab $p_1,q_1\to0$, $q_2\to1$, $w\to1$ that defeated earlier campaigns
(\S\ref{sec:record}): there $\Lb$ collapses while $\Lh$ stays positive.

\subsection{A linear diagonal bound}
\label{sec:rh}

The certificate's ratio rule needs a certified upper bound on $\rho$ that is
\emph{linear and piecewise dyadic-friendly}: over any box, the maximum of
$2\min(x,1-x)$ is attained at a corner, so
$\sqrt{\rh}\le\sqrt{2\min(x,1-x)}$ gives an aggregate bound with no interval
blow-up.

\begin{theorem}[$\rh$ inequality; \texttt{lemma\_rh\_proof.py},
\texttt{lemma\_rh.py}]
\label{thm:rh}
Let $h(u)=-u\log_2u-(1-u)\log_2(1-u)$ on $[0,1]$, with
$h(0)=h(1)=0$, and define
\[
\rh(x)=2(1-x)h(x)-h\bigl((1-x)^2\bigr).
\]
For every $x\in[0,1]$,
\begin{equation}
\rh(x)\le2\min(x,1-x).
\label{eq:rhbound}
\end{equation}
Moreover $\rh(x)/x\to2$ as $x\downarrow0$.
\end{theorem}

\begin{proof}
For $x\in[\tfrac12,1]$, $\rh(x)=2(1-x)h(x)-h((1-x)^2)\le 2(1-x)$, as
$h\le1$ and $h\ge0$. This is the trivial side.

For $x\in[0,\tfrac12]$ the relation $\rh\le2x$ is
$f(x):=h\bigl((1-x)^2\bigr)-2(1-x)\,h(x)+2x\ge0$. The machine-checked
proof (\texttt{lemma\_rh\_proof.py}) is built on the SymPy-exact identity
\begin{equation}
\label{eq:rhexact}
\ln2\cdot f(x)=x^2\ln\frac{2}{x}-x(2-x)\,\ln\Bigl(1-\frac{x}{2}\Bigr),
\qquad 0<x\le\tfrac12,
\end{equation}
with $f(0)=0$ exactly. For $0<x\le\tfrac12$ both summands on the right are
strictly positive ($\ln(2/x)\ge\ln4>0$; $\ln(1-x/2)<0$ negated), so
$f(x)>0$ on $(0,\tfrac12]$: positivity is immediate, with no cancellation. The
near-zero structure is quantified by the explicit remainder
$\ln2\,f(x)=x^2[-\ln x+\ln2+1]+R(x)$ with $-x^3/2\le R(x)\le
-(15/62)\,x^3$, giving $\ln2\,f(x)\ge(\ln32+\tfrac{31}{32})\,x^2\approx
4.434\,x^2$ on $(0,\tfrac1{16}]$; in particular
$f(x)=x^2\log_2(1/x)+O(x^2)=o(x)$, which is the tightness claim. On
$[\tfrac1{16},\tfrac12]$ an adaptive Arb interval cover certifies $f>0$ cell by
cell; the machine proof executes a 44-subinterval cover with minimum certified lower
bound $2.4\times10^{-3}$, after an exact recursive split of the interval.
The \PROVED{} gate checks the SymPy identity, series through order $x^6$,
leading behavior, algebraic factorizations, and remainder bounds; the adaptive
Arb cover checks exact contiguity and positive lower bounds.  The executable
prints \texttt{PROOF COMPLETE: machine-checked} only after both gates pass.
\end{proof}

A companion monotonicity fact supplies the endpoint-safe caps used by
\texttt{cert2.rho\_upper}:

\begin{equation}
\frac{d}{dz}\frac{z}{h(z)}=\frac{h(z)-z\,h'(z)}{h(z)^2}
=\frac{-\log_2(1-z)}{h(z)^2}>0\qquad(0<z<1),
\label{eq:zoverh}
\end{equation}
so $z/h(z)$ increases and by symmetry $h(z)=h(1-z)$ the ratio
$(1-z)/h(z)$ decreases.  These signs are proved symbolically by
\path{cert2._RHO_NUMERATOR_IDENTITY} and asserted at import by
\path{prove_rho_endpoint_monotonicity}.

\section{The certificate}
\label{sec:certificate}

\subsection{The reduced problem and exact arithmetic}
By Corollary~\ref{cor:5param} it remains to certify, for
$\alpha=0.0356069$ and $t_{\mathrm{cert}}=0.3820660112501052$, that
$\Phi_\alpha(\nu)\ge0$ for every two-orbit
$\nu=w\,\delta_{(p_1,q_1)}+(1-w)\,\delta_{(p_2,q_2)}$ with
$\mean(\mu_\nu)\le t_{\mathrm{cert}}$. All universal inequalities are
statements in \emph{Arb ball arithmetic} (\texttt{flint}/\texttt{Arb},
Johansson \cite{arb}): enclosures of $h$, $s^{\star}$ (monotone corners),
$\rh$, and $\sqrt{\rh}$; sound min/max extensions; exact shortest-decimal parsing
of float endpoints into Arb; certified covers. Every rule is an independently sound
lower bound or infeasibility proof: \emph{skipping a rule can only send a box to
the split branch, never clear it.} Numerically derived quantities that steer the
search are labeled NUMERICAL and are never used as proof steps.

The exact orbit swap
$(p_1,q_1,p_2,q_2,w)\mapsto(p_2,q_2,p_1,q_1,1-w)$ preserves
$\mean,\Lb,\Cc,\Ee,\Phi_\alpha$ term by term (asserted exactly and checked on
rational endpoints by \texttt{verify\_orbit\_swap}), so the root
$[0,1]^4\times[\tfrac12,1]$ covers the full cube, and it is split into the
eight exact dyadic $w$-slices
$[\tfrac12,\tfrac9{16}],[\tfrac9{16},\tfrac58],\ldots,[\tfrac{15}{16},1]$
($\texttt{root\_w\_lo}=\texttt{1/2}$, $\texttt{root\_w\_hi}=\texttt{1}$
in the manifest). Campaign parameters (v2 manifests): $\alpha=0.0356069$,
$t=0.3820660112501052$, $\texttt{level\_offset}=10^{-4}$, min box width
$10^{-3}$, collar floor $1.25\times10^{-4}$, face floor $10^{-3}$, face box
budget $10^5$, box budget $2\times10^9$, 24\,h per slice (72\,h in the
relaunched campaign I), 160-bit native
covers with an 80-bit working precision, $\texttt{center\_gate}=-0.02$,
$\texttt{centered5\_max\_width}=1/32$, $\texttt{q2pin}=15/16$,
$\texttt{collar\_edge}=0.01$.

\paragraph{$\lambda$-shifts.}
For $\lambda\ge0$, the shifted functional
$\Phi_{\alpha,\lambda}=\Phi_\alpha+\lambda(\mean-t)$ satisfies
$\Phi_{\alpha,\lambda}\le\Phi_\alpha$ on feasible points
($\mean\le t$), so any certified lower bound on the shifted functional over a box is
a valid lower bound on $\Phi_\alpha$ there. The certified family is
\[
\begin{aligned}
\lambda_{\mathrm{seed}}&=1.3682158100251758,\\
\mathcal L&=\{\,0,\ 1,\ 1.626,\ 2.2,\ 0.98\lambda_{\mathrm{seed}},\
\lambda_{\mathrm{seed}},\ 1.02\lambda_{\mathrm{seed}}\,\}.
\end{aligned}
\]
The displayed decimal is an exact rational seed in the manifest, chosen from
a \NUMERICAL{} face-KKT solve (reported residual below $10^{-35}$).
Correctness does not depend on the seed satisfying any KKT equation.

\subsection{The rule chain}
\texttt{cert3.certify} runs the following cheap-first cascade on a stack of
boxes; each rule writes one trace byte (codes in \S\ref{sec:trace}).

\begin{enumerate}
\item[\texttt{0}] \emph{Mean contractor} (\texttt{diag\_exhaust.mean\_contract});
with exact rationals $l_i=(p_i^{\mathrm{lo}}+q_i^{\mathrm{lo}})/2$ from
shortest-decimal endpoints, every point of the box satisfies
$\mean\ge A(w):=l_2+w(l_1-l_2)$; solve $A(w)\le t$ exactly in
\texttt{Fraction} arithmetic to contract $w$ or prove the box infeasible
(trace code 0).
\item[\texttt{1,2}] \emph{Corner bound} (\texttt{diag\_exhaust.phi\_corner}).
With $m_{\star}=\min(\mean_{\mathrm{hi}},t)$ and
$\kappa_{\star}=2(1-\alpha)(1-m_{\star})-1$ required $>0$, then
$\Phi\ge\kappa_{\star}\,\Lb_{\mathrm{lo}}+\alpha\,\Cc_{\mathrm{lo}}
-(1-\alpha)\,\sigma_{\mathrm{hi}}^2$ over the box (code 1 = infeasible,
code 2 = corner-clear).
\item[\texttt{3}] \emph{Ratio rule} (\texttt{cert2.ratio\_rule}). From the
margin lemma, with a \emph{box-specific} $w$-window
$W=\texttt{hull}(w_{\mathrm{lo}},w_{\mathrm{hi}})$:
$\sigma^2\le\Lb\int\rho\,d\nu\le\Lb\,\rho_{\mathrm{hi}}$, so
\begin{align*}
\Phi&\ge\Lb\bigl[\kappa_t-(1-\alpha)\rho_{\mathrm{hi}}\bigr],\\
\kappa_t&=2(1-\alpha)(1-t)-1
\approx0.191862550011752.
\end{align*}
with $\rho_{\mathrm{hi}}$ the endpoint maximum over $W$ of the convex
combination of the two orbits' certified caps (the combination is linear in $w$;
an unsound slope-form co-variation was identified and replaced by the endpoint
max). This rule is \emph{its own cheapest applicability test} (35.6\,$\mu$s
vs.\ 27.5\,$\mu$s for the corner bound): a hand-tuned geometric pre-filter
(``every atom $\le0.05$ or $\ge0.42$'') delegating applicability was removed
because it fell silent exactly on the near-total-sink slab that matters; see
\S\ref{sec:record}.
\item[\texttt{4--7}] \emph{Centered (KKT-shifted) rules.}
The primary rule is \path{centered5_best}, with
\path{centered5_mixed}, \path{centered5_mixed_swap}, and
\path{centered_w_best} as block-mixed, orbit-swap, and weight-only
fallbacks.  They are gated to widths
$\le1/32$ and require a corner lower bound at least $-0.02$.  Each proves
the lower bound
$\Phi_\lambda\ge\Phi_\lambda(\mathrm{mid})-\sum_i
\sup_{\mathrm{box}}\lvert\partial_i\Phi_\lambda\rvert\cdot w_i$ built from
certified interval gradients over the $\lambda$-family, falling back to
block-mixed forms when the full gradient is unusable at a sink-touching vestigial
orbit, and to the exact orbit-swapped form or the derivative-free weight-only form
on collar-touching boxes.
\item[\texttt{8}] \emph{$q_2$-pin face rule}. When $q_2^{\mathrm{lo}}\ge
\texttt{q2pin}=15/16$, $\sstar(p_2,q_2)=q_2$ is pinned and
$\partial\Phi_\lambda/\partial q_2\le0$ on the strip, so the $q_2$-face
carries the minimum; a four-dimensional face branch-and-bound
(\texttt{face\_bb}) certifies that face with the pinned $\lambda$-subfamily.
\item[\texttt{9}] \emph{Residual}. All five coordinates at their floors (width
$\le10^{-3}$, or $\le1.25\times10^{-4}$ for collar-touching boxes): an
uncleared box, recorded with its center and full widths. A COMPLETE slice has
zero residuals.
\item[\texttt{16+$j$}] \emph{Split of coordinate $j$}. The splitter chooses
the coordinate of largest \emph{influence-weighted} width (width $\times$
influence, where orbit~1's atoms carry at most $w^{\mathrm{hi}}$ and orbit~2's
at most $1-w^{\mathrm{lo}}$; \texttt{split\_by\_influence}, with
\texttt{split\_by\_width} kept beside it behind the \texttt{SPLIT\_CHOICE}
hook). Any split is sound (children cover the parent exactly); the choice affects
only cost.
\end{enumerate}

A box is \emph{terminal} when a rule proves infeasibility, certifies
$\Phi\ge0$, or cannot be split. The verdict is \texttt{COMPLETE} iff the stack
is empty, residuals are zero, and the box and time budgets were not consumed.

\subsection{Proof traces and independent replay}
\label{sec:trace}
Every processed raw DFS node writes exactly one byte (\texttt{cert3-trace-v1}):
codes 0--9 name the terminal rule, code $16+j$ the split coordinate. The trace
size therefore \emph{equals} the processed tally by construction.

\texttt{cert3\_replay.replay\_trace} reconstructs every box from the exact dyadic
root by re-applying the mean contractor and the recorded splits (strict,
representable midpoints required), re-proves \emph{only the claimed rule} per
terminal in Arb interval arithmetic (splits are intrinsically sound), rejects any
residual event, unknown code, trailing event, or non-empty final stack, and
compares all trace-derived tallies against the committed record. Worker fidelity is
enforced by import order: same modules, 160-bit covers, then 80-bit working
precision; the $\lambda$-family is recomputed and checked against the manifest. A
replay failure is a rejection problem.

\begin{itemize}
\item \emph{Adversarial suite}: all eight attack vectors are verified rejections ---
flipped rule byte; injected residual; clear-replaced-by-split; truncation;
extension; unknown code; false infeasibility claim; wrong expected tallies --- each
with the correct diagnostic. An INCOMPLETE slice trace replays mathematically and is
rejected precisely for its pending stack.
\item \emph{Determinism}: campaigns E and F ran slice~0 five hours apart on
different snapshots; both processed exactly 16{,}231{,}621 boxes and their trace
files are byte-identical (SHA-256 \texttt{9eed0cc1382f5485\ldots}). The proof
object is reproducible, not merely re-derivable.
\end{itemize}

\subsection{The immutable campaign protocol (v2)}
\label{sec:protocol}
\path{cert3_par.py} \texttt{create} writes a fresh UUID-named campaign directory.
The seven executable sources are \path{cert3.py}, \path{cert2.py},
\path{cert3_par.py}, \path{diag_exhaust.py}, \path{bound_kkt.py},
\path{arbcore.py}, and \path{entropy.py}.  The eight review sources are
\path{bridge_uc.py}, \path{cert3_collect.py}, \path{cert3_replay.py},
\path{decomposition.py}, \path{lemma_rh_proof.py},
\path{margin_lemma.py}, \path{reduction.py}, and \path{thmB3_proof.py}.
The protocol copies exactly those 15 files into \path{snapshot/}, hashes each
with SHA-256, and atomically writes \path{launch.json} with exact dyadic roots,
canonical lowercase-hex UUID4 run identifiers, and exact parameters.  Workers
run from the snapshot under \texttt{python -B}; the exact listing prevents
imports from preferring stray bytecode or extension modules.

Commit is atomic and append-only. The worker writes its trace to a run-unique
hidden temp file with flush+fsync, records its SHA-256 and size (size must equal
processed), and the supervisor re-hashes the temp before installing it next to the
result via Darwin \texttt{renamex\_np} with \texttt{RENAME\_EXCL}: nothing is
ever deleted or overwritten, and two writers cannot collide.

\path{cert3_collect.py} accepts a campaign only by \emph{re-proving it}.  Its
hardcoded \path{EXPECTED_EXECUTABLE_FILES} (7) and
\path{EXPECTED_REVIEW_FILES} (8) override the manifest's inventory;
the snapshot is re-listed, re-hashed, and must be pristine; roots must be the
exact dyadic partition, run ids canonical UUID4, exit code 0, verdicts
COMPLETE, positive work, and
$\texttt{stack}=\texttt{residual}=\texttt{budget\_boxes}=\texttt{budget\_time}=0$;
trace digests and sizes consistent; then it re-imports the frozen modules from the
snapshot (import location asserted) and \emph{replays all eight traces} before
printing \texttt{COMPOSITE CERTIFICATE}. Hand-authored complete JSON is
insufficient by construction.

\subsection{The campaign record: E, F, G, H, I, and J status}
\label{sec:record}
The certified rational decimal
$t_{\mathrm{cert}}=0.3820660112501052\ge\psi+10^{-4}$ is itself proven
exactly: for $r\in[0,\tfrac32]$, $r\ge\psi\iff r^2-3r+1\le0$, checked
in rationals (\texttt{cert3\_par.target\_relation\_holds}).

\begin{table}[h]
\centering
\footnotesize
\begin{tabular}{r r r r r r r}
\hline
slice & E proc.\ & F proc.\ & G proc.\ & G ratio & H proc. & I proc.\\
\hline
0 & 16,231,621 & 16,231,621 & 15,528,149 & 192,277 & 21,135,611 & 21,135,611$^{\ddagger}$\\
1 & 16,320,717 & 16,320,717 & 15,226,435 & 228,386 & 20,827,065 & 20,827,065$^{\ddagger}$\\
2 & 17,058,935 & 17,058,935 & 15,381,807 & 295,807 & 23,486,265 & 23,486,265$^{\ddagger}$\\
3 & 19,213,165 & 19,213,165 & 17,831,917 & 281,389 & 34,186,855 & 34,186,855$^{\ddagger}$\\
4 & 19,341,141 & 19,341,141 & 19,030,155 & 82,944 & 54,161,482$^{\dagger}$ & 71,405,987$^{\star}$\\
5 & 22,972,035 & 22,972,035 & 22,640,283 & 78,913 & 59,867,104$^{\dagger}$ & 117,709,435$^{\star}$\\
6 & 21,467,895 & 33,614,993 & 29,889,721 & 271,760 & 74,343,239$^{\dagger}$ & 121,456,775$^{\star}$\\
7 & 26,095,385 & 51,360,695 & 52,086,331 & 1,193,438 & 78,257,861 & 78,257,861$^{\ddagger}$\\
\hline
\end{tabular}
\caption{Processed boxes and (for G) ratio clears per $w$-slice in campaigns
E, F, G, H, and I; residuals are $0$ except F slice~7 ($1{,}447{,}470$ sink
residuals), F slices~6--7 and G slice~7 stopped by the wall (stacks 12/18,
24), G slice~7 at $52{,}086{,}331$ boxes with residual~0. E: 2026-08-15,
8\,h/slice, slices 0--5 COMPLETE and replayed. F: 2026-08-16 05:38,
12\,h/slice, slices 0--6 COMPLETE and replayed. G: 2026-08-16 17:56,
24\,h/slice, slices 0--6 COMPLETE and replayed, all rows summing to the
committed 187{,}614{,}798 processed nodes. H (final): 2026-08-17 21:20,
24\,h/slice, slices 0--3 and 7 COMPLETE and replayed, rows summing to the
committed 366{,}265{,}482 processed boxes with residual~0 throughout;
$\dagger$~marks the INCOMPLETE slices 4--6, stopped by the 24\,h wall with
stacks 25/27/30 and residual~0 (time-limited only). I (relaunch): all eight
slice records are committed, rows summing to $488{,}465{,}854$ boxes.
$\ddagger$~marks ``trace SHA-256 identical to campaign H'' (slices 0--3, 7);
$\star$~marks ``I total exceeds H wall'' (slices 4--6). The frozen collector
replayed all eight traces in its own gate and printed \texttt{COMPOSITE
CERTIFICATE} on 2026-08-21. On every replayed slice the binary-tree invariant
$\text{terminals}=\text{splits}+1$ holds exactly.}
\label{tab:campaigns}
\end{table}

\paragraph{How the residual barrier was found and removed (E,F,G).}
Campaign E's slice 6--7 incompleteness was purely wall-clock (12 and 18 pending
stack boxes at the 8\,h wall, zero residuals). Campaign F's slice~7 exposed a
genuinely mathematical-looking wall: 1{,}447{,}470 residuals in 51.4\,M
boxes, all \texttt{sink}, pinned to the razor-thin near-total-sink slab
$p_1,q_1\in[0,10^{-4}]$, $q_2\in[0.9844,1]$, $w\in[0.9920,1]$,
$p_2\in[0.4063,0.4200]$ --- the degenerate family the margin lemma was built
for; the ratio rule cleared 400/400 sampled leaves but was \emph{never called}
there, as a hand-tuned filter demanded every atom $\le0.05$ or $\ge0.42$ and
$p_2$ missed $0.42$ by $1.4\times10^{-2}$. Removing the filter is sound (a
skipped bound only forces a split), the rule is its own cheapest test, and the
extracted slab certifies COMPLETE in 36{,}511 boxes. Campaign G ran without it:
trees \emph{shrunk} (15.2--29.9\,M vs.\ 16.2--33.6\,M) while ratio clears
rose (3.7$\times$ on slice~0); slice~7 passed F's residual onset at box
26.1\,M at residual~0 and reached 52{,}086{,}331 boxes before the 24\,h wall.

Campaign H's only change is the splitter.  The \NUMERICAL{} cost probe
classified slice-7 splits as distributed (48.5\,\% \texttt{other},
35.1\,\% \texttt{mean-boundary}, 16.4\,\% \texttt{sink}), so the remaining
cost appeared computational rather than a new inequality;
\texttt{split\_by\_influence} weights a coordinate's width by its orbit's
weight cap (on slice~7, orbit~2's coordinates are discounted by up to
16$\times$).  The measured changes --- the former residual slab
$36{,}511\to131$ boxes, the tight-obstruction root $231\to109$ boxes, and
$2.8\times$ throughput on $w\in[0.99,1]$ --- are all \NUMERICAL{}.
Campaign H's mathematical acceptance is by replay alone.

\paragraph{Campaign H outcome; campaign I.}
H finished on 2026-08-18: slices 0--3 and 7 COMPLETE --- 21{,}135{,}611,
20{,}827{,}065, 23{,}486{,}265, 34{,}186{,}855 and 78{,}257{,}861 boxes,
stacks and budget flags all zero --- and slices 4--6 INCOMPLETE at the wall
only, at 54{,}161{,}482, 59{,}867{,}104 and 74{,}343{,}239 processed boxes
with stacks 25/27/30. Slice~7, which had defeated E (budget), F (residuals)
and G (budget), finished with 7\,h of headroom; its 78{,}257{,}861-event
trace replayed independently in 2\,h\,11\,m (PASS), as did the four smaller
echoes, every tally matching the committed records and
$\text{terminals}=\text{splits}+1$ exact. Across all 366{,}265{,}482 boxes
the residual count never left zero: the three remaining slices wanted hours,
not mathematics. Campaign I (2026-08-18 21:26\,UTC, workers
\texttt{certI0--7}) committed all eight slices under the same code hash and
exact target with a 72\,h budget per slice. Its totals are
21{,}135{,}611, 20{,}827{,}065, 23{,}486{,}265, 34{,}186{,}855,
71{,}405{,}987, 117{,}709{,}435, 121{,}456{,}775 and 78{,}257{,}861,
for 488{,}465{,}854 boxes overall. The I traces on slices 0--3 and 7 are
byte-identical to H by SHA-256; slices 4--6 have totals exceeding H's wall
totals.  The frozen collector replayed all eight traces and printed
\texttt{COMPOSITE CERTIFICATE} on 2026-08-21: campaign I certifies the
rational target
$t_{\mathrm{cert}}=0.3820660112501052\ge\psi+10^{-4}$.

\paragraph{Second campaign (IN PROGRESS; not a certificate).}
Campaign J, directory
\path{cert3_20260820T181524Z_08ad738a1cf24a08bfc4189e65f9ef44_57a8908ba7b8}
with code hash \texttt{57a8908ba7b8}, is running with a 72-hour budget per
slice at the rational target
$t_J=0.3821660112501052\ge\psi+2\cdot10^{-4}$.  As of 2026-08-21, six slices
are committed \texttt{COMPLETE} with residual zero: slice~0 processed
$21{,}891{,}091$ boxes; slice~1, $21{,}637{,}061$; slice~2,
$24{,}563{,}611$; slice~3, $37{,}877{,}105$; slice~4, $85{,}793{,}971$;
and slice~7, $93{,}281{,}643$.  Slices~5 and 6 are still running.  The frozen
collector verdict has not yet been printed; in particular,
\texttt{COMPOSITE CERTIFICATE} has not been printed for campaign J.
Therefore campaign J remains IN PROGRESS, and this paper makes no
certification claim at $t_J$.

\section{A proof of Liu's positive-semidefiniteness hypothesis}
\label{sec:liu-psd}

Liu's conditionally iid argument \cite[Theorem~13 and \S V]{liu} has two
separate conditional inputs.  Its Hypothesis~1, in \S V-A, is a
positive-semidefiniteness assertion; Hypothesis~2, in \S V-B, asserts the
structure of a nine-parameter global optimizer.  Assuming both, Liu obtains
the union-closed constant
\[
0.382709087918741>c^\star=0.3823455333667027,
\]
strictly beyond the ceiling of the two-strategy method used by our
certificates.  The result below proves Hypothesis~1.  Hypothesis~2 is not
addressed and remains open.  Liu's numerical check of Hypothesis~1 used a
float64 eigenvalue computation on a $2499$-point grid and reported a largest
eigenvalue $1.6311\times10^{-14}$, below the roundoff scale of a dense
computation of that size.

For this section put
\[
a_{\mathrm L}(s)=s(1-s),\qquad
z(s,t)=(1-s)(1-t)+a_{\mathrm L}(s)a_{\mathrm L}(t).
\]
Liu's hypothesis says that
\begin{equation}
\iint h(z(s,t))\,d\mu(s)\,d\mu(t)\le0
\label{eq:liu-hypothesis}
\end{equation}
for every finite signed measure $\mu$ annihilating
\[
1,\qquad s,\qquad a_{\mathrm L}(s)=s(1-s).
\]
These three functions span all polynomials of degree at most two.  Thus the
hypothesis has codimension three, as does Liu's projector
$[1,i/n,a_{\mathrm L}(i/n)]$; the positive-semidefiniteness terminology refers
to the negative of the entropy kernel on this subspace.

\subsection{Exact reduction to an unprojected kernel}

Set
\[
u=1-s,\qquad v=1-t,\qquad A=uv,\qquad B=st,
\]
so that $a_{\mathrm L}(s)a_{\mathrm L}(t)=AB$ and
$z=A(1+B)$.  With
\[
G(x)=x+(1-x)\ln(1-x)
     =\sum_{k=2}^{\infty}\frac{x^k}{k(k-1)}
\]
on $[0,1]$, the entropy has the exact expansion
\begin{equation}
\begin{aligned}
\ln2\,h(z)
 &=-z\ln u-z\ln v-z\ln(1+B)+z-G(z)\\
 &=-z\ln u-z\ln v-z\ln(1+B)+z
   -\sum_{k=2}^{\infty}\frac{z^k}{k(k-1)}.
\end{aligned}
\label{eq:liu-expansion}
\end{equation}
Products involving an endpoint logarithm are understood by continuous
extension.  The three displayed terms $-z\ln u$, $-z\ln v$, and $+z$ are
exactly the pieces discarded by Liu's projection.  Indeed,
\[
\begin{aligned}
-z\ln u&=-(u\ln u)v-(us\ln u)(vt),\\
-z\ln v&=-u(v\ln v)-(us)(vt\ln v),\\
z&=uv+(us)(vt),
\end{aligned}
\]
and every rank-one summand has, in at least one variable, a polynomial factor
of degree at most two.  This is precisely why all three constraints
$1,s,a_{\mathrm L}(s)$ are needed.

The rank-one kernel $AB=a_{\mathrm L}(s)a_{\mathrm L}(t)$ is itself
annihilated.  Moving that null summand from $z=A+AB$ into the remainder rewrites
\eqref{eq:liu-expansion} as
\[
\ln2\,h(z)=-z\ln u-z\ln v+A+R(s,t),
\]
where
\begin{equation}
R(s,t)
=A\bigl[B-(1+B)\ln(1+B)\bigr]
 -\bigl[z+(1-z)\ln(1-z)\bigr].
\label{eq:liu-R}
\end{equation}
Consequently, for every $\mu$ satisfying Liu's three constraints,
\begin{equation}
\ln2\iint h(z(s,t))\,d\mu(s)\,d\mu(t)
=\iint R(s,t)\,d\mu(s)\,d\mu(t).
\label{eq:liu-reduction}
\end{equation}
Thus it suffices to prove the stronger statement that $R$ is negative
semidefinite without any projection.

\begin{theorem}[Liu's positive-semidefiniteness hypothesis]
\label{thm:liu-psd}
The continuous kernel $R$ in \eqref{eq:liu-R} is negative semidefinite on
$[0,1]$: for every $\nu\in\MM([0,1])$,
\[
\iint R(s,t)\,d\nu(s)\,d\nu(t)\le0.
\]
In particular, its integral operator is negative semidefinite on
$L^2[0,1]$, and \eqref{eq:liu-hypothesis} holds.
\end{theorem}

\begin{proof}
Continue to regard $A=uv$ and $B=st$ as kernels.  Extend
$G(x)=x+(1-x)\ln(1-x)$ to $[-1,1]$.  Since
$G(-B)=-B+(1+B)\ln(1+B)$, define
\[
\phi(B)=G\bigl(A(1+B)\bigr)+A\,G(-B)=-R.
\]
Treating $A$ as fixed, direct symbolic differentiation gives
\begin{align*}
\phi(0)&=G(A),&
\phi'(0)&=-A\ln(1-A),\\
\phi''(B)&=\frac{A}{(1+B)\,[1-A(1+B)]}.&&
\end{align*}
Taylor's theorem with integral remainder therefore gives the exact identity
\begin{equation}
\begin{aligned}
-R
={}&G(A)-AB\ln(1-A)\\
&\quad+B^2\int_0^1(1-r)
\frac{A}{(1+rB)\,[1-A(1+rB)]}\,dr.
\end{aligned}
\label{eq:liu-taylor}
\end{equation}

The first two kernels in \eqref{eq:liu-taylor} are positive semidefinite by
nonnegative power series and the Schur product theorem:
\[
G(A)=\sum_{k=2}^{\infty}\frac{A^k}{k(k-1)},\qquad
-AB\ln(1-A)=\sum_{k=1}^{\infty}\frac{A^{k+1}B}{k}.
\]
Indeed, $A^k=u^kv^k$ and
$A^{k+1}B=(u^{k+1}s)(v^{k+1}t)$ are rank-one positive-semidefinite kernels.

It remains to treat the integral remainder.  For $0\le r\le1$ write
\[
D_r(s,t)=(1+rB)\,[1-A(1+rB)].
\]
Let $\mathbf m(x)=(1,x,x^2,x^3)^{\mathsf T}$.  In this polynomial basis the
kernel has coefficient matrix
\begin{equation}
\begin{aligned}
M(r)&=
\begin{pmatrix}
0&1&0&0\\
1&-r-1&2r&0\\
0&2r&-r(r+2)&r^2\\
0&0&r^2&-r^2
\end{pmatrix},\\
D_r(s,t)&=\mathbf m(s)^{\mathsf T}M(r)\mathbf m(t).
\end{aligned}
\label{eq:liu-M}
\end{equation}
Exact coefficient extraction gives
\[
\det M(r)=-2r^3,
\]
and the coefficients of $\det(\lambda I-M(r))$, in descending powers of
$\lambda$, are
\begin{equation}
1,\quad (r+1)(2r+1),\quad 4r^3+2r-1,\quad
-2r(r^3-r^2+r+1),\quad -2r^3.
\label{eq:liu-charpoly}
\end{equation}
For $0<r\le1$ their signs are
$(+,+,\mathord\pm,-,-)$, with a zero coefficient omitted.  There is exactly
one Descartes sign change.  Since $M(r)$ is symmetric and its determinant is
nonzero, it has inertia $(1\text{ positive},3\text{ negative})$.

For completeness, an exact $LDL^{\mathsf T}$ factorization makes the resulting
Lorentz decomposition explicit.  Put
\[
\begin{aligned}
\zeta_0(s)&=-\frac{2rs^2-(r+1)s+1}{r+1},
&\zeta_1(s)&=1+2rs^2,\\
\zeta_2(s)&=\frac{s^2(r+2-rs)}{r+2},
&\zeta_3(s)&=s^3,
\end{aligned}
\]
and
\[
\begin{aligned}
\alpha_r(s)&=\frac{\zeta_1(s)}{\sqrt{r+1}},\\
\mathbf b_r(s)&=\left(
\sqrt{r+1}\,\zeta_0(s),\
\sqrt{r(r+2)}\,\zeta_2(s),\
r\sqrt{\frac{2}{r+2}}\,\zeta_3(s)
\right).
\end{aligned}
\]
Then, also at $r=0$ by continuity,
\begin{equation}
D_r(s,t)=\alpha_r(s)\alpha_r(t)
-\langle\mathbf b_r(s),\mathbf b_r(t)\rangle.
\label{eq:liu-lorentz}
\end{equation}
The positive coordinate cannot degenerate away from the corner.  On the
diagonal,
\begin{align}
D_r(s,s)
&=(1+rs^2)\,[1-(1-s)^2(1+rs^2)]\notag\\
&\ge s(2-2s+2s^2-s^3)>0
\qquad(0<s\le1).
\label{eq:liu-diagonal}
\end{align}
For the inequality, apply concavity of
$q\mapsto q[1-(1-s)^2q]$ on $1\le q\le1+s^2$ and check its two endpoints.
The cubic $p(s)=2-2s+2s^2-s^3$ has
$p'(s)=-3s^2+4s-2$, whose discriminant is $-8$; hence $p$ decreases from
$2$ to $1$ on $[0,1]$.

Set $\mathbf c_r(s)=\mathbf b_r(s)/\alpha_r(s)$.  Equation
\eqref{eq:liu-lorentz} and \eqref{eq:liu-diagonal} give
$\|\mathbf c_r(s)\|<1$ for $s>0$.  Thus Cauchy--Schwarz gives absolute
convergence of the Gram series
\begin{equation}
\frac1{D_r(s,t)}
=\frac1{\alpha_r(s)\alpha_r(t)}
\sum_{k=0}^{\infty}
\left\langle
\mathbf c_r(s)^{\otimes k},\mathbf c_r(t)^{\otimes k}
\right\rangle
\qquad(s,t>0).
\label{eq:liu-geometric}
\end{equation}
Every summand is a positive-semidefinite kernel.  Multiplication by the
rank-one kernel
\[
AB^2=(1-s)s^2(1-t)t^2
\]
therefore shows that $AB^2/D_r$ is positive semidefinite on $(0,1]^2$.
It extends continuously by zero to the closed square: near $(0,0)$,
$D_r(s,t)=s+t+O((s+t)^2)$ uniformly for $r\in[0,1]$, whereas
$AB^2=O(s^2t^2)$.  Positive semidefiniteness at boundary points follows by
taking limits of the finite Gram matrices.  Integration against the
nonnegative weight $(1-r)\,dr$ preserves positive semidefiniteness.

The middle term of \eqref{eq:liu-taylor} also has the asserted continuous
extension at the only singular corner.  If $m=\max(s,t)$, then
$1-A=s+t-st\ge m$ and $B=st\le m^2$, so
\[
0\le AB[-\ln(1-A)]\le m^2(-\ln m)\longrightarrow0.
\]
Hence every term in \eqref{eq:liu-taylor} extends to $[0,1]^2$ as a
positive-semidefinite kernel.  Their sum is $-R$, so $R$ is negative
semidefinite.  The signed-measure statement follows either from the displayed
Gram representations or by approximating against the continuous kernels, and
\eqref{eq:liu-reduction} proves Liu's Hypothesis~1.
\end{proof}

\paragraph{Scope.}
Theorem~\ref{thm:liu-psd} settles Liu's Hypothesis~1 only.  Hypothesis~2, the
nine-parameter global-optimization assertion in Liu's \S V-B, remains open.
Consequently the value $0.382709087918741$ remains a conditional
union-closed constant; this paper makes no unconditional claim at that value.

\begin{remark}[Why the hypothesis resisted, and tightness]
Three natural proof routes can be ruled out.  First, termwise
Cauchy--Schwarz/telescoping domination cannot close: the negative double-series
family is confined to the index band $i\le k$, while every positive term is of
the form $us^j$ with $j\ge3$, so the available negative weight at the first
positive index is exactly zero.  Second, a uniform spectral-gap comparison
cannot close because the negative-side operator is compact on an
infinite-dimensional space and has lower spectral edge zero.  Third, a Schur
factorization through $(1+\theta st)^{-1}$ is invalid because that kernel is
not positive semidefinite.  Already at $\theta=1/2$ and points $1/4,3/4$, its
$2\times2$ Gram determinant is the exact negative rational
$-131072/1657425$; larger \NUMERICAL{} Gram tests had minimum eigenvalues from
$-1.02$ at $\theta=0.25$ to $-9.47$ at $\theta=1$ for $60$ and $200$ nodes.
The inequality is tight at the corner: the resolved \NUMERICAL{} generalized
ratio is $0.9999999906$, and the factorization identifies the exact source of
the lost margin as $D_r(0,0)=0$.
\end{remark}

\paragraph{Corroboration only (not part of the proof).}
The symbolic identities and diagnostics are implemented in
\texttt{liu\_kernel.py}, \texttt{liu\_R\_nsd.py}, \texttt{liu\_tail.py}, and
\texttt{liu\_moment\_cert.py}.  On a $300$-node grid,
\eqref{eq:liu-taylor} reproduced $-R$ to $8.9\times10^{-16}$; the three pieces
in its displayed order had minimum eigenvalues
$-1.09\times10^{-14}$, $-1.92\times10^{-15}$, and
$-4.78\times10^{-16}$.  The sampled kernels $A/D_r$ had minimum eigenvalue
about $-7\times10^{-14}$ for
$r\in\{0,0.1,0.25,0.5,0.75,1\}$.  Independently, exact moments followed by
Arb arithmetic certified the following upper bounds for compressions of $R$
to polynomials of degree at most $D$, in the orthonormal shifted-Legendre basis
$\sqrt{2n+1}\,P_n(2s-1)$:
\[
\begin{aligned}
D=8:\quad
&\lambda_{\max}\le-1.64335523965845934982130\times10^{-10},\\
D=12:\quad
&\lambda_{\max}\le-1.39301806562092415191600\times10^{-14},\\
D=16:\quad
&\lambda_{\max}\le-1.37614421207490684661199\times10^{-18}.
\end{aligned}
\]
These are strictly negative validated enclosures on the stated finite
subspaces.  They corroborate, but do not replace, the continuum proof above.

\appendix
\section{Reproducibility}
\label{sec:verify}

\subsection{Frozen artifact and collector command}
The certified campaign directory is
\begin{quote}\small
\path{cert3_20260818T212601Z_425f109c15b64a6198785c6cebbbdaab_2f23a58ebdb8}.
\end{quote}
From the repository's \texttt{math/} directory, the complete replay is the
following single command:
\begin{quote}
\footnotesize
\texttt{.venv/bin/python -B }\path{uc/campaigns/cert3_20260818T212601Z_425f109c15b64a6198785c6cebbbdaab_2f23a58ebdb8/snapshot/cert3_collect.py}\texttt{ }\path{uc/campaigns/cert3_20260818T212601Z_425f109c15b64a6198785c6cebbbdaab_2f23a58ebdb8/launch.json}
\end{quote}
Exit code 0 and the final line \texttt{COMPOSITE CERTIFICATE} are jointly
required.  The accepted collector has SHA-256
\begin{quote}\small
\path{95d09322a7821f41850ef1ce77345e5e93eb4fa570cc0e7f58f817716cf681ec}.
\end{quote}
The full evidence bundle used for independent transfer is
\begin{quote}\small
\path{/Users/jinleic/jinleic-workspace/data/uc-certificate-psi-plus-1e-4-2026-08-21-v2.zip}.
\end{quote}

\subsection{Pinned environment}
The launch manifest pins CPython 3.14.3, flint 0.9.0, mpmath 1.3.0,
numpy 2.5.2, scipy 1.18.0, and sympy 1.14.0 on macOS arm64.  Workers and the
collector run the frozen sources with \texttt{python -B}.  The accepted code
SHA-256 is
\path{2f23a58ebdb8b14275284e2ffefca7862373f7baf5d354a06e607cff8d78ce05}.
Never import the snapshot modules from another program: importing can create a
bytecode cache before the module disables bytecode writes, thereby changing
the immutable snapshot listing.

\subsection{Committed tallies and replay time}
Every campaign-I slice was committed with
\[
\texttt{stack}=\texttt{residual}=\texttt{budget\_boxes}
=\texttt{budget\_time}=0.
\]
The collector record is:
\begin{center}
\small
\begin{tabular}{r r r}
\hline
slice & processed boxes & replay elapsed (s)\\
\hline
0 & 21{,}135{,}611 & 2{,}436.6\\
1 & 20{,}827{,}065 & 2{,}442.1\\
2 & 23{,}486{,}265 & 2{,}721.5\\
3 & 34{,}186{,}855 & 4{,}326.3\\
4 & 71{,}405{,}987 & 9{,}746.8\\
5 & 117{,}709{,}435 & 25{,}329.8\\
6 & 121{,}456{,}775 & 21{,}622.3\\
7 & 78{,}257{,}861 & 9{,}115.2\\
\hline
total & 488{,}465{,}854 & 77{,}740.6\\
\hline
\end{tabular}
\end{center}
Thus the practical expected serial replay wall time on the pinned host is
about $77{,}741$ seconds, or 21 hours 36 minutes.  Hardware and system load can
change elapsed time but cannot change any accepted tally or verdict.

\subsection{Acceptance checks and residual trust base}
The collector performs the protocol of \S\ref{sec:protocol}: it checks its
hardcoded source inventories, pristine snapshot listing and hashes, exact
dyadic roots, unique UUID4 run identifiers, worker exit codes, COMPLETE
verdicts, positive work, zero stack/residual/budget fields, and trace
digest/size consistency.  It then imports only the frozen modules and
mathematically replays every claimed discharge in Arb arithmetic.  A
hand-authored result JSON cannot pass this gate.

The residual trust base is explicit: CPython semantics and startup behavior;
the pinned flint/Arb implementation; mpmath, numpy, scipy, and sympy at the
versions above; Darwin \texttt{renamex\_np} with
\texttt{RENAME\_EXCL}; SHA-256; the hashed Python sources; and the
\CITED{} results used in the analytic reduction --- Riesz,
Banach--Alaoglu, and Stone--Weierstrass as presented by van Neerven
\cite{vanneerven}, Bauer's principle \cite{bauer,brupedra}, the moment-set
theorem \cite{winkler,pinelis}, and Cambie's single entropy-to-union-closed
implication \cite[Question~2 and \S4]{cambie}.

\paragraph{Artifact map.}
\texttt{bridge\_uc.py} checks the $s^\star$ identity and folding bridge;
\texttt{decomposition.py} checks Theorem~\ref{thm:A};
\texttt{reduction.py} checks Theorem~\ref{thm:Ap},
Corollary~\ref{cor:Fid}, Corollary~\ref{cor:psi}, and
Lemma~\ref{lem:acrit}; \texttt{thmB3\_proof.py} checks
Theorem~\ref{thm:B3}; \texttt{margin\_lemma.py} checks
Lemma~\ref{lem:margin}; \texttt{lemma\_rh\_proof.py},
\texttt{cert2.py}, and \texttt{arbcore.py} check
Theorem~\ref{thm:rh} and its caps; \texttt{cert3.py} writes the certificate
and traces; \texttt{cert3\_replay.py} replays them;
\texttt{cert3\_par.py} enforces the campaign protocol; and
\texttt{cert3\_collect.py} performs final acceptance.
\texttt{liu\_kernel.py} checks the codimension-three reduction;
\texttt{liu\_R\_nsd.py} checks the Taylor and Lorentz identities and the Arb
compressions; and \texttt{liu\_tail.py} and
\texttt{liu\_moment\_cert.py} audit the failed routes and corroborating
calculations.

\section*{Acknowledgments}
This note is part of the union-closed certification project.  The companion
proof ledger and status documents are maintained with the repository
artifacts.  All certified computations reported here were performed on the
pinned macOS arm64 environment described in Appendix~\ref{sec:verify}.

\bibliographystyle{plain}
\bibliography{refs}

\end{document}
